\chapter{Literature Review}\label{ch:lit}
PaperMind draws on four lines of work: analysis and classification of examination questions, automatic question
generation, large language models with retrieval and tool use, and label-efficient learning. This chapter reviews each
and identifies the gap the project addresses.

\section{Related Work}
\subsection{Representing and classifying examination questions}
Classical information retrieval represents a text as a weighted bag of terms. Sparck Jones introduced inverse document
frequency to down-weight common terms~\cite{sparckjones1972}, and Salton and Buckley systematised TF-IDF weighting
schemes~\cite{salton1988}; with cosine similarity these remain strong baselines for short technical texts whose meaning
is carried by domain vocabulary. Bloom's taxonomy, revised by Anderson and Krathwohl~\cite{anderson2001}, organises
cognitive demand into six levels (Remember, Understand, Apply, Analyze, Evaluate, Create) and is widely used to audit
question papers. Mohammed and Omar classified examination questions into Bloom levels using a modified TF-IDF combined
with word2vec features~\cite{mohammed2020}, showing that lexical cues---chiefly the directive verb---carry most of the
signal. Contextual encoders such as BERT~\cite{devlin2019} and sentence-level variants such as
Sentence-BERT~\cite{reimers2019} and MiniLM~\cite{wang2020minilm} produce dense embeddings that enable zero-shot matching
of a question against textual topic descriptions.

\subsection{Automatic question generation}
Early question generation was rule- and template-based: Mitkov and Ha generated multiple-choice items from
textbooks using term extraction and distractor selection~\cite{mitkov2003}, and Heilman and Smith over-generated
questions by syntactic transformation and ranked them statistically~\cite{heilman2010}. Du \emph{et al.}\ framed
question generation as sequence-to-sequence learning with attention~\cite{du2017}. Kurdi \emph{et al.}\ surveyed
93 educational question generation systems and found that most generate isolated factual questions from a text passage,
few control difficulty or cognitive level, and evaluation of answer correctness is rare~\cite{kurdi2020}. None of these
systems model the pattern of a whole examination.

\subsection{Large language models, retrieval and tool use}
The Transformer~\cite{vaswani2017} enabled LLMs that perform new tasks from instructions and a few
examples~\cite{brown2020}; Google's Gemini family adds native multimodal input~\cite{gemini2023}. LLMs are, however,
prone to hallucination---fluent but unfaithful output~\cite{ji2023}. Retrieval-augmented generation (RAG) conditions
generation on retrieved documents to ground it~\cite{lewis2020}. Self-consistency improves reasoning accuracy by sampling
several solutions and checking agreement~\cite{wang2023sc}, an idea that motivates the blind re-solving of answer keys
in this project. ReAct~\cite{yao2023} and Toolformer~\cite{schick2023} showed that LLMs can decide when to call external
tools, the basis of tool-using assistants such as the one built here.

\subsection{Label-efficient learning}
Knowledge distillation trains a small student model to imitate a larger teacher~\cite{hinton2015}. Weak supervision
replaces hand labels with noisy programmatic sources~\cite{ratner2017}. More recently LLMs themselves have been used as
annotators: Wang \emph{et al.}\ showed GPT-3 labels can cut labelling cost substantially~\cite{wang2021gpt3label}, and
Gilardi \emph{et al.}\ found LLM annotations rival crowd workers on several text-classification
tasks~\cite{gilardi2023}. PaperMind uses an LLM as teacher and a TF-IDF + logistic-regression student~\cite{pedregosa2011}
so that the deployed system needs no deep-learning runtime.

\subsection{Document digitisation}
Optical character recognition (OCR) engines such as Tesseract~\cite{smith2007} convert scanned pages to text; modern
platform OCR and multimodal LLMs extend this to layout-rich documents. Mathematical notation remains a known weak point
of text-layer extraction, which this project confirms on real question papers.

\begin{table}[t]
\centering
\caption{Summary of related work and comparison with PaperMind}
\label{tab:related}
\small
\begin{tabularx}{\textwidth}{@{}>{\raggedright\arraybackslash}p{3.3cm}>{\raggedright\arraybackslash}Xccc@{}}
\toprule
\textbf{Work} & \textbf{Approach} & \textbf{Pattern} & \textbf{Key check} & \textbf{Personal}\\
\midrule
Mitkov \& Ha~\cite{mitkov2003} & Template MCQ generation from text & -- & -- & --\\
Heilman \& Smith~\cite{heilman2010} & Syntactic over-generate and rank & -- & -- & --\\
Du \emph{et al.}~\cite{du2017} & Neural seq2seq question generation & -- & -- & --\\
Mohammed and Omar~\cite{mohammed2020} & Bloom classification (TF-IDF + word2vec) & partial & -- & --\\
Lewis \emph{et al.}~\cite{lewis2020} & Retrieval-augmented generation & -- & -- & --\\
Wang \emph{et al.}~\cite{wang2023sc} & Self-consistency reasoning & -- & \checkmark & --\\
Yao \emph{et al.}~\cite{yao2023} & Tool-using LLM agents & -- & -- & --\\
\textbf{PaperMind} & Analysis + RAG generation + distilled labels + agent & \checkmark & \checkmark & \checkmark\\
\bottomrule
\end{tabularx}
\end{table}

\section{Research Gap}
Table~\ref{tab:related} summarises the landscape. Question analysis work classifies questions but does not turn the
analysis into generation; question generation work produces individual questions from passages but ignores the
structure of a real examination; and LLM-based generation is ungrounded and unchecked by default. No reviewed system
(i) learns an examination's subject, topic and format distribution from its own past papers, (ii) uses that
distribution to constrain an LLM that generates complete papers, (iii) verifies the generated answer keys, and
(iv) closes the loop with a student's measured performance. PaperMind is designed to fill exactly this gap.

\section{Chapter Summary}
This chapter reviewed lexical and dense representations of questions, Bloom's taxonomy classification, the evolution
of question generation from templates to LLMs, retrieval augmentation, self-consistency and tool use, and LLM-assisted
labelling. The combination of these ideas into a single pattern-aware, verified and personalised system is the
contribution of this work.
